\section{Введение: битва архитектур и четырёхслойный налог}
\label{sec:intro}

\subsection{Соглашение фон Неймана и узкое место, которое оно не убрало}

Компьютер с хранимой программой, описанный в черновике EDVAC фон Неймана
1945 года, поместил инструкции и данные в одну память, адресуемую одним
каналом. Его победа была не заявкой на оптимальность, а на \emph{универсальность}:
одна память, один последовательный поток, любая программа. Всякая машина
общего назначения с тех пор --- его потомок. Но универсальность куплена
структурной ценой, которую Джон Бэкус назвал в Тьюринговской лекции 1977 года
--- \emph{узким местом фон Неймана}~\cite{backus1977}: канал процессор--память,
через который всё вычисление приходится прокачивать «по слову за раз». Семьдесят
лет архитектуры были, в значительной части, кампанией против этого канала:
кэши, предвыборка, внеочередное исполнение, широкие векторы и глубокие иерархии
памяти --- всё это способы прятать его, а не убрать. Причина, по которой его
нельзя убрать, теперь --- доминирующий факт области. После конца деннардовского
масштабирования (ок.\ 2006) связывающее ограничение --- энергия, а энергия во
фон-неймановской машине доминируется не арифметикой, а \emph{перемещением
данных}: операция с плавающей точкой стоит порядка пикоджоуля, тогда как
выборка её операндов из внекристальной памяти стоит на один-два порядка
больше. Мы считаем дёшево и двигаем данные дорого, а архитектура построена
вокруг неправильного из двух.

\subsection{Четыре подсистемы, спроектированные порознь}

Под узким местом лежит более глубокий и менее обсуждаемый факт. Современная
вычислительная система --- и эмфатически \emph{обучающаяся} --- построена из
четырёх подсистем, которые задумываются, оптимизируются и верифицируются
\emph{независимо}, а затем принуждаются сосуществовать:

\begin{enumerate}[leftmargin=1.7em,itemsep=1pt]
  \item \textbf{интерконнект} --- сеть-на-кристалле или топология фабрики,
        переносящая состояние между вычислительными узлами;
  \item \textbf{код коррекции ошибок} (ECC) --- отдельный слой избыточности,
        привинченный к памяти и связям, чтобы пережить отказы;
  \item \textbf{набор инструкций} (ISA) --- алгебра примитивных операций,
        которые машина умеет выполнять;
  \item \textbf{правило адаптации} --- для обучающейся машины, процедура
        тренировки/оптимизации, меняющая её параметры.
\end{enumerate}

У каждой своя литература, свои оптимумы и свои проектировщики. Поскольку они
проектируются порознь, их интерфейсы приходится согласовывать, и всякое
согласование --- \emph{налог рассогласования}, оплачиваемый трижды:

\begin{table}[H]
\centering
\small
\caption{Четыре слоя проектирования суть отдельные подсистемы в каждой
существующей машине, каждая несёт независимую цену. \TALOS{} платит каждую
цену однажды, в одной структуре.}
\label{tab:four-layers}
\begin{tabular}{@{}p{2.4cm}p{4.6cm}p{5.2cm}@{}}
\toprule
\textbf{Слой} & \textbf{Как делается сегодня} & \textbf{Какой налог берёт} \\
\midrule
Интерконнект & Кроссбар / mesh / NoC, ради полосы & Площадь, статич.\ мощность, сложность маршрутизации; all-to-all растёт как $\binom{N}{2}$ \\
ECC & Хэмминг/BCH/LDPC привинчены к хранению и связям; surface-коды для кубитов & Выделенные проверочные биты ($+12{,}5\%$ для ECC-DRAM) или $\sim\!10^3\times$ физич./логич.\ кубита \\
ISA & Фиксированный набор опкодов, со-эволюция с компилятором & Семантический разрыв; «стиль фон Неймана», осуждённый Бэкусом \\
Адаптация & Отдельный тренировочный цикл (обратное распространение) над замороженной ISA & Память под активации/градиенты; раздел train/infer; энергия $\propto$ перемещению данных \\
\bottomrule
\end{tabular}
\end{table}

\subsection{Претенденты и что каждый из них исправил}

Альтернативные архитектуры были каждая атакой на \emph{один} из этих четырёх
слоёв, никогда на все четыре:

\begin{itemize}[leftmargin=1.5em,itemsep=1pt]
  \item \textbf{Гарвардская} разделила память команд и данных --- микро-оптимизация
        границы ISA/интерконнект, выжившая в кэшах и DSP.
  \item \textbf{Dataflow}~\cite{dennis1975} убрал счётчик команд: операции
        срабатывают по готовности операндов. Он реформировал слой
        \emph{инструкции} (управление данными, не потоком), но не предложил
        принципиальной топологии, кода или правила обучения.
  \item \textbf{Систолические массивы}~\cite{kung1982} навязали ритмичный,
        локальный, ближайшесоседский \emph{интерконнект} --- и выиграли
        матричные блоки сегодняшних GPU и TPU --- но топология была инженерным
        выбором под одно ядро, не выведенной структурой.
  \item \textbf{Нейроморфные} машины~\cite{mead1990} (TrueNorth, Loihi,
        SpiNNaker) со-локализовали память и счёт и сделали субстрат событийным
        и диссипативным --- реформировав интерконнект \emph{и} адаптацию сразу
        --- но с \emph{ad hoc} топологией, без родной коррекции ошибок и без
        принципиального объекта состояния.
  \item \textbf{Вычисления в памяти} и мемристорные кроссбары выполняли
        умножение-накопление \emph{в} массиве памяти, атакуя узкое место прямо,
        но не несли ни динамики, ни самомодели.
  \item \textbf{Квантовые} машины сделали слой инструкции унитарным (обратимым,
        значит дешёвым по Ландауэру) --- и затем вынуждены были привинчивать
        коррекцию ошибок с разорительными накладными, потому что код и
        вычисление суть разные структуры, декогерирующие друг против друга.
\end{itemize}

Каждая из них была подлинным продвижением по одной оси. Ни одна не объединила
четыре, потому что ни у одной не было \emph{принципа}, из которого можно было
бы \emph{вывести} единую структуру. Топология всегда была догадкой; код всегда
был отдельным; правило обучения всегда жило над ISA. Область искала структуру,
которая одновременно хорошая сеть, хороший код, хорошая алгебра инструкций и
хорошее правило обучения --- и трактовала это как четыре поиска.

\subsection{Тезис: архитектурный монизм}

Этот документ сообщает, что у четырёх поисков один ответ и что ответ
\emph{вынужден}, а не выбран. Унитарный Голономный Монизм (УГМ) --- математическая
теория, выводящая структуру жизнеспособной самомоделирующей системы из
единственного примитива, матрицы когерентности
$\Gam\in\Dens(\mathbb{C}^7)$~\cite{uhm2026} --- доказывает, что на семимодовом
субстрате интерконнект, код коррекции ошибок, алгебра инструкций и правило
обучения суть \emph{не четыре объекта, а один}: плоскость Фано $\Fano$. Мы
называем это \emph{архитектурным монизмом}. Раздел~\ref{sec:collapse} излагает
это как две теоремы:

\begin{itemize}[leftmargin=1.5em,itemsep=1pt]
  \item \textbf{Теорема~\ref{thm:collapse} (Коллапс).} Семь прямых $\Fano$
        суть одновременно интерконнект; та же инцидентность --- проверочная
        геометрия кода Хэмминга $[7,4,3]$; октонионные структурные константы на
        этих прямых суть умножение алгебры инструкций; а $G_2$-ковариантная
        неподвижная точка этой алгебры --- правило обучения. Все четыре суть
        функции одной инцидентной структуры $7\times7$. Налог рассогласования
        тождественно ноль, потому что нет интерфейсов для согласования.
  \item \textbf{Теорема~\ref{thm:unique} (Уникальность).} Этот коллапс
        реализуем при \emph{единственном} числе узлов, $N=7$. Совпадение
        «порядок проективной плоскости $=$ длина Хэмминга» повторяется (при $N =
        7, 31, \dots$), но лишь при $N=7$ алгебра инструкций несёт ещё и
        нетривиальный ассоциатор --- третьепорядковую структуру, на которой
        держится весь проект. Неассоциативность разрывает ничью в пользу семи и
        нигде более.
\end{itemize}

Архитектура, которая из этого следует, \TALOS{},\footnote{По имени бронзового
автомата греческого мифа, первого искусственного существа --- здесь физического
тела искусственного разума. Имя --- кодовое, не аббревиатура.} поэтому не есть
предложение среди предложений. Это аппаратная форма математического объекта.
Его состояние --- матрица плотности в $784$~байта, помещающаяся в L1-кэш; его
единственная инструкция --- УГМ-линдбладиан $\LindOmega$; его сеть --- семь шин;
его отказоустойчивость --- эта сеть; его обучение --- третий член инструкции; а
его потребление отслеживает то, как быстро оно учится, а не то, как тяжело оно
считает.

\subsection{Честные границы, указанные первыми}

Две теоремы УГМ ограничивают то, что \TALOS{} может и не может утверждать, и мы
выносим обе вперёд, чтобы ничто ниже не читалось сверх того, что оно есть.

\begin{itemize}[leftmargin=1.5em,itemsep=1pt]
  \item \textbf{Субстратная мультипликативность (T-153).} Теорема субстратной
        замкнутости УГМ доказывает, что вычисление субстрат-\emph{независимо}:
        выбор физического вложения --- чистая $G_2$-калибровка. Нет единственной
        физической машины --- есть единственная \emph{логическая} архитектура и
        семейство реализующих её субстратов. \TALOS{} --- импеданс-согласованный
        член этого семейства (\S\ref{sec:realizability}), не единственная
        возможность. Это сила (переносимость), а не слабость.
  \item \textbf{Вычислительная граница (BQP).} УГМ доказывает, что
        регенеративный член $\Regen$ не даёт \emph{никакого} ускорения за
        пределы класса сложности BQP. \TALOS{} --- не машина быстрее квантовой.
        Его преимущество --- в энергии, адаптивности и отказоустойчивости ---
        способность-на-ватт и родное обучение --- никогда в классе сложности.
        Мы повторяем это везде, где читатель мог бы перечитать заявление сверх
        меры.
\end{itemize}

Остаток документа --- спецификация. Раздел~\ref{sec:collapse} доказывает две
теоремы. Разделы~\ref{sec:substrate}--\ref{sec:composition} специфицируют
состояние, фабрику, код-как-топологию, энергетическую модель, самомодель и закон
композиции. Раздел~\ref{sec:realizability} даёт лестницу реализуемости от
сегодняшней GPU-эмуляции до управляемо-диссипативного аналогового устройства.
Раздел~\ref{sec:abi} специфицирует двоичный интерфейс приложений к \SYNARC{},
программному разуму, для которого \TALOS{} --- тело. Раздел~\ref{sec:redteam}
подвергает проект явной red-team, а раздел~\ref{sec:benchmarks} сводит его к
фальсифицируемым инженерным измерениям против фон-неймановской базы. Каждое
число машинно-проверено.
